% Generated draft: requires mathematical and source review.
\documentclass[11pt]{article}
\usepackage[margin=25mm]{geometry}
\usepackage{fontspec}
\setmainfont{Latin Modern Roman}
\usepackage{amsmath,amssymb,mathtools,unicode-math}
\setmathfont{Latin Modern Math}
\usepackage{enumitem,needspace,xurl,hyperref}
\setlength{\parindent}{0pt}
\setlength{\parskip}{5pt}
\setlength{\emergencystretch}{3em}
\newcommand{\N}{\mathbb{N}}
\newcommand{\Z}{\mathbb{Z}}
\newcommand{\Q}{\mathbb{Q}}
\newcommand{\R}{\mathbb{R}}
\newcommand{\C}{\mathbb{C}}
\newcommand{\F}{\mathbb{F}}
\newcommand{\A}{\mathbb{A}}
\newcommand{\PP}{\mathbb{P}}
\newcommand{\E}{\mathbb{E}}
\newcommand{\eps}{\varepsilon}
\newcommand{\dd}{\,\mathrm{d}}
\begin{document}
\section*{UnsolvedMath ID: 30006587}
\textbf{Source problem code:} catalogue entry
\subsection{Definitions and mathematical statement}
Let \(A=\Q\langle z_2,z_3,\ldots\rangle\) be the free associative unital \(\Q\)-algebra on the indicated generators. Its basis consists of finite words in these generators, including the empty word. The supplied record denotes the intended multiple Eisenstein series associated with \(w\) by \(G(w)\), with dependence on a variable \(\tau\).

The record describes an operator \(\theta\), called a derivation in its background assessment, by
\[
\theta(w)=-D\bigl(\varphi(w)\bigr),
\]
and says that its construction uses the operator \(\mathrm{Drop1}\). It supplies no definitions of \(D\), \(\varphi\), or \(\mathrm{Drop1}\), nor the spaces on which they act. Thus this formula does not by itself define \(\theta\) explicitly.

\textbf{Conjectural target, subject to the unresolved definitions below.} With the intended multiple Eisenstein evaluation and operators, establish
\[
\forall w\in A,\qquad
2\pi i\,\frac{d}{d\tau}G(w)(\tau)
   =G\bigl(\theta(w)\bigr)(\tau)
\]
at every \(\tau\) in the intended common domain. The source record does not specify that domain, the normalization or regularization of \(G\), its evaluation on operator outputs, or its convention on the empty word. These conventions must be recovered before this is a fully specified mathematical assertion. The wording also calls \(w\) a word while quantifying over the whole algebra; extending a wordwise identity to all algebra elements requires an appropriate linearity convention that is not explicitly supplied.

The original extracted text additionally reports that \(R(z_2,z_2)\) yields the weight-six relation
\[
6G_{3,3}-3G_{4,2}-G_6=0.
\]
This is retained as a supplied illustrative relation, not as an independently verified assertion or a definition of the target. Neither \(R\) nor the conventions for the indexed series are defined in the record. The original extraction is incomplete around the conjecture itself, so the displayed derivative identity above follows the supplied reconstructed statement.
\subsection{Short English statement}
The supplied conjecture says that differentiation of a multiple Eisenstein series is represented by the operator \(\theta=-D\circ\varphi\), with the factor \(2\pi i\) shown above. Definitions of the series and operators are missing, so this remains a statement-only draft requiring human mathematical review.
\subsection{Sources}
\begin{itemize}
\item Local UnsolvedMath registry, supplied source record: permanent ID \(30006587\), exact source problem code \texttt{catalogue entry}, title \emph{Derivative Formula for Multiple Eisenstein Series}, in the Oberwolfach Reports Open Problems collection; proposed year \(2026\). No individual proposer is identified.
\item \emph{Multiple Zeta Values and Modular Forms} (2026), Oberwolfach Reports, report 7/2026, pp.\ 433--486, DOI 10.4171/owr/2026/7. Supplied source and literature link: \url{https://doi.org/10.4171/owr/2026/7}.
\item Bachmann--Kanno, \emph{Relations and Derivatives of Multiple Eisenstein Series}. Supplied literature link: \url{https://arxiv.org/abs/2602.08176}. The original extracted conjecture refers to a reference labelled [2], but the supplied record does not explicitly identify that label with a bibliography entry.
\item Literature assessment dated \texttt{2026-08-22}: supplied dated curation, not independently verified. It labels the conjecture open and states that the 2026 primary paper proves structural results but does not prove the full assertion involving \(\theta\) and the formal multiple-Eisenstein relations. The background explicitly requests independent verification of time-sensitive status and cited claims.
\item The registry labels the statement \texttt{corrected\_verified} and reports that it was reconstructed and checked against the official Oberwolfach report. This is supplied provenance, not verification performed for this draft. The original extracted statement is visibly incomplete.
\end{itemize}
\end{document}
